% Fichier généré automatiquement par tools/generate_corrections.py.
% Source : STAT/STAT-03-Demarche/STOCK/02_R_02/02_R_02.tex
% Statut : partiel (2/3 question(s) renseignée(s), 2 reprise(s)).
\begin{corrigebox}[Corrigé partiel extrait de la banque]
\CorrigeQuestion{1}
\CorrigeACompleter{Aucun élément de réponse n'est présent dans la source d'origine pour cette question.}
\CorrigeQuestion{2}
\CorrigeSource{Réponse reprise d'un exercice homonyme de la banque ; la provenance exacte est indiquée dans le commentaire du fichier.}
% Réponse reprise de : CIN/CIN-02-VitesseAcceleration/02_R_02/02_R_02.tex
On a directement $\torseurcin{V}{1}{0} = \torseurl{\dot{\theta}\vk{0}}{R\dot{\theta}\vj{1}}{B}$.
\CorrigeQuestion{3}
\CorrigeSource{Réponse reprise d'un exercice homonyme de la banque ; la provenance exacte est indiquée dans le commentaire du fichier.}
% Réponse reprise de : CIN/CIN-02-VitesseAcceleration/02_R_02/02_R_02.tex
$\vectg{B}{1}{0} = \deriv{\vectv{B}{1}{0}}{\rep{0}}=R\ddot{\theta}\vj{1}-R\dot{\theta}^2\vi{1}$.
(En effet,  $\deriv{\vj{1}}{\rep{0}} = \deriv{\vj{1}}{\rep{1}} + \vecto{1}{0}\wedge \vj{1} $ $=\vect{0}+\dot{\theta}\vk{0}\wedge \vj{1}$ $=-\dot{\theta}\vi{1}$.)
\end{corrigebox}
